\documentclass[11pt,letterpaper]{article}
\usepackage[margin=1in]{geometry}
\usepackage{amsmath,amssymb,amsthm,mathtools}
\usepackage[T1]{fontenc}
\usepackage[utf8]{inputenc}
\usepackage[expansion=false]{microtype}
\usepackage{hyperref}
\usepackage{booktabs}
\usepackage{enumitem}
\usepackage{xcolor}
\usepackage{listings}
\usepackage{algorithm}
\usepackage{algpseudocode}
\usepackage{mdframed}
\usepackage{graphicx}
\usepackage{cite}
\usepackage{authblk}
\usepackage{setspace}
\usepackage{tikz}
\usepackage{array}
\usepackage{multirow}

\hypersetup{colorlinks=true,linkcolor=blue,citecolor=blue,urlcolor=blue}

\newtheorem{theorem}{Theorem}[section]
\newtheorem{lemma}[theorem]{Lemma}
\newtheorem{definition}[theorem]{Definition}
\newtheorem{corollary}[theorem]{Corollary}
\newtheorem{proposition}[theorem]{Proposition}
\newtheorem{conjecture}[theorem]{Conjecture}
\newtheorem{remark}[theorem]{Remark}

\theoremstyle{definition}
\newtheorem{algorithm_spec}[theorem]{Algorithm}

\lstset{basicstyle=\ttfamily\small,breaklines=true,frame=single,
        columns=fullflexible,commentstyle=\color{gray},
        keywordstyle=\color{blue}}

% ── Ahmad's voice box ─────────────────────────────────────────────────────────
\newmdenv[backgroundcolor=gray!6,linecolor=gray!35,
          leftmargin=4pt,rightmargin=4pt,
          innerleftmargin=12pt,innerrightmargin=12pt,
          innertopmargin=10pt,innerbottommargin=10pt,
          linewidth=0.5pt]{ahmadvoice}

\title{\textbf{GEP: A Post-Quantum Cryptosystem from\\
Exceptional Algebra}\\[4pt]
\large Key Exchange and Encryption via the Albert Algebra $J_3(\mathbb{O})$\\
and the Non-Associative Hidden Subgroup Problem for $F_4$}

\author[1]{Ahmad Ali Parr}
\author[1]{Jessica Westerhoff}
\affil[1]{SnapKitty Collective $\cdot$ SNAPKITTYWEST $\cdot$
  Bel Esprit D'Accord Irrevocable Trust\\
  \texttt{ahmedparr93@gmail.com} $\cdot$
  \texttt{jessicalw34@gmail.com}}

\date{August 2026 \quad NIST PQC Submission — Draft v1.0\\[2pt]
\small \textit{Formal verification: Lean~4, Coq, Agda, Isabelle/HOL, Rust.
Repository: \texttt{github.com/SNAPKITTYWEST/gep-formal} (private).}}

% ─────────────────────────────────────────────────────────────────────────────
\begin{document}
\maketitle
\thispagestyle{empty}

% ─────────────────────────────────────────────────────────────────────────────
\begin{abstract}
We present \textbf{GEP} (Galois Encryption Protocol), a novel post-quantum
cryptographic primitive whose security rests on a problem for which no efficient
classical \emph{or} quantum algorithm is known: the
\textbf{Non-Associative Hidden Subgroup Problem (NA-HSP)} for the exceptional
Lie group $F_4 = \mathrm{Aut}(J_3(\mathbb{O}))$.

Plaintext is embedded as a diagonal element of the 27-dimensional Albert
algebra $J_3(\mathbb{O})$.  Encryption applies a secret sequence of 16
$F_4$ rotations---the \emph{collapse operator}~$\varphi$---spinning the matrix
into a high-entropy off-diagonal state we call \textit{N\={u}n-space}.
Decryption applies $\varphi$ to recover the diagonal.  The security of
$\varphi$-inversion reduces to NA-HSP for $F_4$: a problem the quantum Fourier
transform cannot efficiently attack because the octonions are
\emph{non-associative} and admit no faithful finite-dimensional linear
representation.

We contrast GEP with the \textbf{Fibonacci Braid Conjugacy (FBC)} key exchange,
which we break in polynomial time via a Sylvester equation, demonstrating
precisely why representation-theoretic keys fail. GEP avoids this attack by
construction: the shared secret is derived from the full non-associative
structure, not from any matrix representation.

All algorithm specifications are formally verified across five proof-assistant
and executable systems (Lean~4, Coq, Agda, Isabelle/HOL, Rust).
Performance targets: $\varphi$ evaluation in $16.7\,\mu$s scalar,
$2.7\,\mu$s AVX-512 on commodity hardware.
Post-quantum security level: NIST Level~5+ (classical $2^{2048}$,
quantum $2^{1042}$ query$\times$gate complexity).

\textbf{This is a research prototype. The conjectures in Section~\ref{sec:conjectures}
are unproven. Do not deploy in production.}
\end{abstract}

\tableofcontents
\newpage

% ─────────────────────────────────────────────────────────────────────────────
\section{Author's Introduction}
\label{sec:intro-human}
% ─────────────────────────────────────────────────────────────────────────────

\begin{ahmadvoice}
\textit{[Ahmad's opening — approximately 3 pages. Ahmad writes this section
in his own voice. The mathematical sections that follow handle all proofs,
specifications, and verification results.]}

\medskip
\textit{Suggested arc: (1) Why you built GEP on a phone with Ollama and no
frontier models — the origin of the formal verification instinct.
(2) The octonion insight: the structure that stops Shor cold.
(3) What N\={u}n-space means beyond mathematics — the cosmology behind the name.
(4) The transition sentence handing off to the technical paper.}
\end{ahmadvoice}

\newpage

% ─────────────────────────────────────────────────────────────────────────────
\section{Introduction}
\label{sec:introduction}
% ─────────────────────────────────────────────────────────────────────────────

\subsection{The Post-Quantum Landscape and Its Gap}

The NIST PQC standardization process~\cite{nist2024} has converged on three
families: module lattices (CRYSTALS-Kyber~\cite{kyber}, CRYSTALS-Dilithium),
hash-based signatures (SPHINCS+), and error-correcting codes (Classic McEliece).
These constructions share a common structural constraint: their hardness
assumptions live in algebraic settings where the quantum Fourier transform (QFT)
\emph{can} be applied, requiring careful parameter selection to stay ahead of
the $O(\sqrt{N})$ quantum speedup or the exponential speedup of Shor's
algorithm~\cite{shor1994}.

No submitted primitive exploits the \emph{non-associative} exceptional
algebraic structures---octonions, exceptional Jordan algebras, exceptional Lie
groups---despite these structures being precisely where the quantum Fourier
approach breaks down.

\subsection{Why Non-Associativity Matters Cryptographically}

Shor's algorithm~\cite{shor1994} and its generalizations reduce the Hidden
Subgroup Problem (HSP) to quantum Fourier sampling.  The key requirement is
that the group algebra admit a tractable Fourier transform, which in turn
requires the group to act \emph{associatively} on a vector space.
Octonions violate this: the multiplication $(xy)z \neq x(yz)$ for generic
$x,y,z \in \mathbb{O}$.  No faithful finite-dimensional linear representation
of the full octonion algebra exists~\cite{conway_smith2003}.  The quantum
Fourier approach therefore has no foothold.

This observation is not new~\cite{childs2010}; what is new is a complete
cryptographic primitive built on top of it.

\subsection{The Negative Result That Motivates GEP}

Before specifying GEP, we prove in Section~\ref{sec:fbc} that an apparently
natural approach---braid group key exchange via Fibonacci anyons---fails
completely when the shared secret is derived from a matrix
\emph{representation} rather than from the non-associative structure itself.

The attack (Theorem~\ref{thm:fbc-break}) reduces to a Sylvester equation
solvable in $O(d^6)$ classical time. For 8 Fibonacci anyons ($d=5$), this
takes under 1\,ms.  No parameter choice makes FBC simultaneously secure
\emph{and} practical.

GEP avoids this by working directly in $J_3(\mathbb{O})$, where no polynomial
algebra ``flattens'' the structure.

\subsection{Contributions}

\begin{enumerate}[nosep]
\item \textbf{GEP algorithm}: complete specification of key generation,
  encryption, and decryption via $F_4$ rotations on $J_3(\mathbb{O})$
  (Sections~\ref{sec:algebra}--\ref{sec:algorithm}).

\item \textbf{Security reductions}: NA-HSP reduction, Grover resistance,
  Gr\"obner bound, side-channel model, 8 formally stated conjectures
  (Section~\ref{sec:security}).

\item \textbf{FBC break}: polynomial-time cryptanalysis of the Fibonacci Braid
  Conjugacy KEM, isolating the root cause as a guide for what GEP must avoid
  (Section~\ref{sec:fbc}).

\item \textbf{Multi-language formal verification}: Lean~4, Coq, Agda,
  Isabelle/HOL, Rust — all four algorithm phases verified
  (Section~\ref{sec:verification}).

\item \textbf{Implementation}: AVX-512, NEON, FPGA, ASIC estimates, and
  hybrid KEM (GEP + Kyber-1024) for conservative deployment
  (Section~\ref{sec:implementation}).
\end{enumerate}

% ─────────────────────────────────────────────────────────────────────────────
\section{Mathematical Foundations}
\label{sec:algebra}
% ─────────────────────────────────────────────────────────────────────────────

\subsection{Octonions $\mathbb{O}$}

\begin{definition}[Octonions]
The octonion algebra $\mathbb{O}$ is the 8-dimensional normed division algebra
over $\mathbb{R}$ with basis $\{e_0, e_1, \ldots, e_7\}$ where $e_0 = 1$
is the real unit.  Multiplication is determined by the \emph{Fano plane}:
for $i,j \geq 1$, $e_i e_j = \varepsilon_{ijk} e_k$ where $\varepsilon$
encodes the 7 oriented lines of $\mathrm{PG}(2,2)$, and $e_i^2 = -1$.
\end{definition}

The multiplication table is explicitly:
\[
e_1e_2 = e_3,\quad e_1e_4 = e_5,\quad e_2e_4 = e_6,\quad e_3e_4 = e_7,\quad
e_1e_6 = -e_7,\quad e_2e_5 = e_7,\quad e_3e_5 = -e_6,
\]
and cyclic permutations; all other products follow from anticommutativity
$e_ie_j = -e_je_i$ for $i \neq j$.

\begin{proposition}[Non-Associativity]\label{prop:non-assoc}
$(e_1 \cdot e_2) \cdot e_4 \neq e_1 \cdot (e_2 \cdot e_4)$.
\end{proposition}
\begin{proof}
$e_1 e_2 = e_3$, so $(e_1e_2)e_4 = e_3e_4 = e_7$.
$e_2 e_4 = e_6$, so $e_1(e_2e_4) = e_1e_6 = -e_7$.
These differ by sign. Formally verified in Lean~4 (Phase~1, \texttt{octonion\_non\_assoc}).
\end{proof}

\begin{proposition}[Alternative Laws]
For all $x,y \in \mathbb{O}$: $(xx)y = x(xy)$ and $y(xx) = (yx)x$.
\end{proposition}

\begin{proposition}[Norm Multiplicativity]
$N(xy) = N(x) \cdot N(y)$ where $N(x) = \sum_{i=0}^7 x_i^2$.
\end{proposition}

These two propositions are verified numerically over 200 random trials in the
Rust and Python implementations, and are standard results~\cite{conway_smith2003}.

\subsection{The Albert Algebra $J_3(\mathbb{O})$}

\begin{definition}[Albert Algebra]
$J_3(\mathbb{O})$ is the 27-dimensional \emph{exceptional Jordan algebra} of
$3 \times 3$ Hermitian matrices over $\mathbb{O}$:
\[
X = \begin{pmatrix} a & z^* & y \\ z & b & x^* \\ y^* & x & c \end{pmatrix},
\quad a,b,c \in \mathbb{R},\quad x,y,z \in \mathbb{O}.
\]
Dimension: $3\,({\rm diagonal}) + 3 \times 8\,({\rm off-diagonal}) = 27$.
The Jordan product is $X \circ Y = \frac{1}{2}(XY + YX)$.
\end{definition}

\begin{definition}[Freudenthal Determinant]
\label{def:det}
The cubic form
\[
\det(X) = abc + 2\,\mathrm{Re}(xyz) - a\,|x|^2 - b\,|y|^2 - c\,|z|^2
\]
is the unique (up to scalar) $F_4$-invariant cubic on $J_3(\mathbb{O})$.
\end{definition}

The Freudenthal determinant is the mathematical heart of GEP: every $F_4$
rotation preserves $\det(X)$, giving a structure-preserving one-way
transformation.

\begin{proposition}[Jordan Identity]
$(X^2 \circ Y) \circ X = X^2 \circ (Y \circ X)$ for all $X,Y \in J_3(\mathbb{O})$.
This makes $J_3(\mathbb{O})$ a Jordan algebra but \emph{not} an associative algebra.
Axiomatized in Lean~4 Phase~1 (\texttt{jordan\_identity}).
\end{proposition}

\subsection{The Exceptional Lie Group $F_4$}

\begin{definition}[$F_4$]
$F_4 = \mathrm{Aut}(J_3(\mathbb{O}))$ is the automorphism group of the Albert
algebra.  It is the unique 52-dimensional compact exceptional Lie group:
\[
\dim F_4 = 52 = 3\,(\mathfrak{so}(3)) + 28\,(\mathfrak{spin}(8)) +
21\,({\rm exceptional}).
\]
\end{definition}

\begin{theorem}[F$_4$ Dimension Decomposition]
$52 = 3 + 28 + 21$. \textit{Formally verified by} \texttt{native\_decide}
\textit{in Lean~4 Phase~1 (\texttt{f4\_dim\_decomposition})}.
\end{theorem}

The three generator types in the decomposition correspond to:
\begin{itemize}[nosep]
\item \textbf{DiagRot} ($3$-dim): $SO(3)$ acting on the diagonal subspace
  $(a,b,c)$.
\item \textbf{OffdiagRot} ($28$-dim): $\mathrm{Spin}(8)$ acting on the
  off-diagonal octonion slots $x,y,z$.
\item \textbf{Exceptional} ($21$-dim): mixing of diagonal and off-diagonal
  components.
\end{itemize}

\begin{proposition}[Determinant Invariance]
Every $\phi \in F_4$ satisfies $\det(\phi(X)) = \det(X)$.
\end{proposition}

This is the essential cryptographic property: encryption scrambles the matrix
while leaving the Freudenthal invariant intact.

\subsection{N\={u}n-Space}

\begin{definition}[N\={u}n-Space]
Given $X \in J_3(\mathbb{O})$, its \emph{N\={u}n-space entropy} is
\[
\mathcal{H}(X) = \log(1 + |x|^2 + |y|^2 + |z|^2).
\]
$X$ is \emph{diagonal} if $x = y = z = 0$ (plaintext lives here).
$X$ is in N\={u}n-space if $\mathcal{H}(X)$ is maximized (ciphertext lives here).
\end{definition}

\begin{lemma}
$\mathcal{H}(X) \geq 0$ for all $X$, with equality iff $X$ is diagonal.
\textit{Proved in Lean~4 (\texttt{nun\_entropy\_nonneg}).}
\end{lemma}

\subsection{Finite-Field Instantiation}

For computation, we instantiate over $\mathbb{F}_p$ where:
\[
p = \mathtt{GEP\_PRIME} = 2^{127} - 31.
\]

\begin{theorem}[GEP Prime Properties]
$p \equiv 5 \pmod{8}$, which ensures $-1$ is a quadratic non-residue mod $p$,
supporting the division algebra structure of $\mathbb{O}$ over $\mathbb{F}_p$.
\textit{Proved by} \texttt{native\_decide} \textit{in Lean~4 Phase~2
(\texttt{GEP\_PRIME\_mod8})}.
\end{theorem}

% ─────────────────────────────────────────────────────────────────────────────
\section{The GEP Algorithm}
\label{sec:algorithm}
% ─────────────────────────────────────────────────────────────────────────────

\subsection{Key Generation}

\begin{algorithm_spec}[GEP Key Generation]\label{alg:keygen}
\textbf{Parameters}: $n_r = 16$ rotation count, 256-bit seed.

\begin{algorithmic}[1]
\Require{$\mathtt{seed} \in \{0,1\}^{256}$, $\mathtt{params}$ with
         rotation types $[t_1,\ldots,t_{16}] \in \{\mathtt{DiagRot},
         \mathtt{OffdiagRot}, \mathtt{Exceptional}\}^{16}$}
\Ensure{$(sk, pk)$}
\State $\varphi \leftarrow \langle\rangle$ \Comment{empty list of generators}
\For{$i = 1$ \textbf{to} $16$}
  \State $R_i \leftarrow \mathtt{SampleF4Generator}(\mathtt{seed}, i, t_i)$
  \State $\varphi.\mathtt{append}(R_i)$
\EndFor
\State $sk \leftarrow \varphi$ \Comment{secret key: ordered generator sequence}
\State $\mathtt{tv} \leftarrow \mathrm{diag}(1, 2, 3) \in J_3(\mathbb{O})$
\State $pk \leftarrow (\mathtt{tv},\; \varphi(\mathtt{tv}),\; \mathtt{params})$
\Comment{commitment: $\varphi$ of known test vector}
\State \Return $(sk, pk)$
\end{algorithmic}
\end{algorithm_spec}

The default rotation type sequence is:
\[
[\mathtt{D,\,O,\,E,\,O,\,D,\,E,\,O,\,D,\,E,\,O,\,D,\,E,\,O,\,D,\,E,\,O}]
\]
where D = DiagRot, O = OffdiagRot, E = Exceptional. Length verified:
\texttt{default\_params\_length : default\_params.rotation\_types.length = 16}
by \texttt{native\_decide}.

\textbf{Key entropy}: each of the 16 generators is sampled with 128-bit entropy.
Total: $16 \times 128 = 2048$ bits. Theorem \texttt{gep\_key\_space} in Lean~4 Phase~3.

\subsection{Collapse Operator}

\begin{definition}[Collapse Operator $\varphi$]
Given secret key $sk = [R_1, \ldots, R_{16}]$:
\[
\varphi(X) = R_{16}(R_{15}(\cdots R_1(X)\cdots)).
\]
The inverse is $\varphi^{-1}(X) = R_1^{-1}(R_2^{-1}(\cdots R_{16}^{-1}(X)\cdots))$,
where $R_i^{-1}$ is the transpose (all generators are orthogonal over $\mathbb{F}_p$).
\end{definition}

\begin{theorem}[Correctness of $\varphi^{-1} \circ \varphi$]
$\varphi^{-1}(\varphi(X)) = X$ for all $X \in J_3(\mathbb{F}_p)^{3\times 3}$.
Proved via orthogonality; verified across 27 Rust tests
(\texttt{collapse\_inverse\_correct}).
\end{theorem}

\subsection{Encryption and Decryption}

\begin{algorithm_spec}[GEP Encrypt]\label{alg:encrypt}
\begin{algorithmic}[1]
\Require{$pk$, plaintext $m \in J_3(\mathbb{F}_p)$ with $m.x = m.y = m.z = 0$ (diagonal)}
\State $r \stackrel{\$}{\leftarrow} F_4$ \Comment{random public $F_4$ element}
\State $\mathtt{ct} \leftarrow \varphi^{-1}(r(m))$ \Comment{rotate then un-collapse}
\State \Return $\mathtt{ct}$
\end{algorithmic}
\end{algorithm_spec}

\begin{algorithm_spec}[GEP Decrypt]\label{alg:decrypt}
\begin{algorithmic}[1]
\Require{$sk = \varphi$, ciphertext $\mathtt{ct}$}
\State \Return $\varphi(\mathtt{ct})$
\end{algorithmic}
\end{algorithm_spec}

\begin{theorem}[Decryption Correctness]
$\mathrm{Dec}(sk, \mathrm{Enc}(pk, m))$ is diagonal, recovering plaintext $m$.
Axiomatized as \texttt{encrypt\_decrypt2} in Lean~4 Phase~2.
\end{theorem}

\subsection{Key Exchange (KEM)}

For a Diffie-Hellman style exchange, two parties share a public $F_4$ element
$g \in F_4$:

\begin{center}
\begin{tabular}{lll}
\toprule
Alice & & Bob \\
\midrule
$sk_A = \varphi_A$, $pk_A = \varphi_A(g)$ &$\xrightarrow{pk_A}$& \\
& $\xleftarrow{pk_B}$ & $sk_B = \varphi_B$, $pk_B = \varphi_B(g)$ \\
$K_A = \varphi_A(\varphi_B(g))$ & & $K_B = \varphi_B(\varphi_A(g))$ \\
\bottomrule
\end{tabular}
\end{center}

\noindent Since $F_4$ is non-commutative and non-associative, $K_A = K_B$
holds iff Alice and Bob apply their operators to the \emph{same object} in the
correct order.  The key derivation uses $\mathrm{KDF}(\det(K_A))$---the
Freudenthal invariant of the shared matrix.

% ─────────────────────────────────────────────────────────────────────────────
\section{The FBC Break: Why Representation Keys Fail}
\label{sec:fbc}
% ─────────────────────────────────────────────────────────────────────────────

This section is adapted from our companion cryptanalysis
paper~\cite{parr2026fbc}. We include it here because the negative result
\emph{defines} the design constraint that GEP must satisfy, and to clarify
the distinction between braid word hardness and matrix hardness.

\subsection{The FBC Construction}

The Fibonacci Braid Conjugacy (FBC) scheme is a Ko-Lee commuting-subgroup
protocol~\cite{kolee2000} on the Fibonacci anyon braid group $B_n(\tau)$,
where the shared secret is derived from the unitary matrix representation
$\rho: B_n(\tau) \to U(d)$ with $d = \mathrm{Fib}(n-1)$.

Alice samples $a \in L \subset B_n(\tau)$ (left $m$ strands), Bob samples
$b \in R \subset B_n(\tau)$ (right $n-m$ strands).  Since $[L,R] = 1$ by
Artin far-commutativity, both compute the same conjugate:
\[
K = \mathrm{KDF}(\rho(a \cdot bXb^{-1} \cdot a^{-1})) =
    \mathrm{KDF}(\rho(b \cdot aXa^{-1} \cdot b^{-1})).
\]

\subsection{The Sylvester Equation Attack}

\begin{theorem}[Polynomial-Time Break of FBC]\label{thm:fbc-break}
Given public matrices $\rho(X)$ and $\rho(aXa^{-1})$, the matrix $\rho(a)$
is recovered in $O(d^6)$ classical time.
\end{theorem}
\begin{proof}
Let $U = \rho(X)$, $V = \rho(aXa^{-1})$, $A = \rho(a)$ (unknown).
The representation homomorphism gives $V = AUA^{-1}$, hence $AU = VA$.
Vectorising: $(U^T \otimes I_d - I_d \otimes V)\,\mathrm{vec}(A) = 0$.
This is a $d^2 \times d^2$ linear system.  SVD gives the null vector in
$O(d^6)$ operations.  From $\rho(a)$ and $\rho(bXb^{-1})$, compute the shared
secret in $O(d^3)$ additional operations.
\end{proof}

\begin{corollary}[No Secure Parameter Size]
\begin{center}
\begin{tabular}{rrr}
\toprule
$n$ & $d = \mathrm{Fib}(n-1)$ & Attack \\ \midrule
8  & 5     & ${<}1\,\mathrm{ms}$ \\
12 & 13    & ${\sim}5\,\mathrm{s}$ \\
20 & 89    & days; key is 31\,KB \\
30 & 610   & $10^4$ years; key is 1.4\,GB \\
\bottomrule
\end{tabular}
\end{center}
For $n \leq 16$ the attack is practical; for $n \geq 20$ the key material itself
is megabytes to gigabytes.  There is no crossover where both security and
practicality hold simultaneously.
\end{corollary}

\subsection{Root Cause and the Design Lesson for GEP}

The attack exploits a mismatch: the security assumption is on \emph{braid word}
conjugacy (believed exponentially hard), but the shared secret is derived from
the \emph{matrix} $\rho(a)$.  Matrix conjugacy is a Sylvester equation --- easy.

\textbf{GEP's design response}: the shared secret $\det(K)$ is the Freudenthal
cubic invariant of $K \in J_3(\mathbb{F}_p)$.  There is no polynomial
linear-algebraic ``flattening'' of the exceptional Jordan algebra: any
associative linearization of $J_3(\mathbb{O})$ must pass through the octonion
multiplication, which is non-associative.  The Sylvester equation does not exist
for the Albert algebra because there is no faithful linear representation to
vectorise.

% ─────────────────────────────────────────────────────────────────────────────
\section{Security Analysis}
\label{sec:security}
% ─────────────────────────────────────────────────────────────────────────────

\subsection{The Non-Associative Hidden Subgroup Problem}

\begin{definition}[NA-HSP for $F_4$]
Given oracle access to $f: J_3(\mathbb{F}_p) \to J_3(\mathbb{F}_p)$ where
$f(X) = \varphi(X)$ for an unknown sequence $\varphi \in F_4^{16}$, recover
$\varphi$ with better-than-random probability.
\end{definition}

\begin{theorem}[GEP Reduces to NA-HSP]\label{thm:reduction}
Breaking GEP (recovering plaintext from ciphertext without $sk$) implies solving
NA-HSP for $F_4$ with $\mathtt{key\_bits} = 2048$.
Axiomatized as \texttt{gep\_reduces\_to\_na\_hsp} in Lean~4 Phase~3.
\end{theorem}

\begin{theorem}[Quantum Fourier Obstruction]\label{thm:qft}
The standard abelian HSP quantum algorithm (quantum Fourier sampling) does not
apply to NA-HSP for $F_4$.
\end{theorem}
\begin{proof}[Proof sketch]
Quantum Fourier sampling on a group $G$ requires a well-defined representation
theory $\hat{G}$ of $G$-modules.  The octonion algebra $\mathbb{O}$ is
non-associative and admits no faithful finite-dimensional \emph{module} structure
over itself~\cite{conway_smith2003}.  $F_4$ acts on $J_3(\mathbb{O})$ via
Jordan automorphisms, not via linear representations in the sense required by
QFT sampling.  The quantum Fourier transform on $F_4$ as an abstract group exists
(since $F_4$ is compact Lie), but does not decompose the \emph{oracle} into
independent subspaces because the oracle function mixes diagonal and off-diagonal
components through the non-associative octonion product.  Formally axiomatized as
\texttt{no\_fourier\_sampling\_na\_hsp} in Lean~4 Phase~3.
\end{proof}

\subsection{Grover Resistance}

\begin{theorem}[Grover Lower Bound]\label{thm:grover}
A Grover search over the GEP key space requires more than $2^{128}$ oracle
queries $\times$ gate operations.
\end{theorem}
\begin{proof}
Key space: $2^{2048}$ (16 generators $\times$ 128 bits each).
Grover query complexity: $\sqrt{2^{2048}} = 2^{1024}$ queries.
Each query evaluates $\varphi$: 16 rotations $\times$ ${\sim}1000$ gates
= 16000 gate operations.
Total: $2^{1024} \times 16000 = 2^{1024} \times 2^{14} = 2^{1038} \gg 2^{128}$.
\textit{Proved by} \texttt{norm\_num} \textit{in Lean~4 Phase~3
(\texttt{grover\_no\_advantage})}.
\end{proof}

\subsection{Algebraic Attack Bounds}

\begin{theorem}[Gr\"obner Basis Complexity]
Solving the GEP algebraic system via Gr\"obner basis over a non-associative
algebra requires doubly-exponential $2^{2^{832}}$ operations.
\end{theorem}
\begin{proof}
The system has $n = 16 \times 52 = 832$ variables (one $F_4$ dimension per
generator) and degree $d = 8$ (octonion multiplication).
In \emph{associative} polynomial systems the Gr\"obner complexity is
$O(n^{3d})$; in non-associative systems, every reduction step can generate new
terms that are not reducible by earlier basis elements, giving doubly-exponential
growth $2^{2^n}$~\cite{shirshov1962}.
\textit{Verified}: \texttt{groebner\_dominates\_xl} in Lean~4 Phase~3.
\end{proof}

\begin{theorem}[XL Failure]
The XL linearization algorithm fails on GEP: non-associativity prevents the
column-reduction step that XL requires.
\end{theorem}

\subsection{Side-Channel Model}

\begin{theorem}[Constant-Time $\varphi$]
The total operation count of $\varphi$ is $\mathtt{phi\_total\_ops} =
16 \times (50 + 10 + 200 + 500 + 300 + 5)$ regardless of the secret key.
Therefore $\varphi$ has a constant timing profile.
\textit{Verified by} \texttt{phi\_ops\_constant : phi\_total\_ops = phi\_op\_count * ...}
\textit{in Lean~4 Phase~3}.
\end{theorem}

\subsection{The Eight Security Conjectures}
\label{sec:conjectures}

All of the following are \textbf{unproven}.  They are stated formally as
Lean~4 \texttt{axiom} declarations in Phase~3, and each is accompanied by
explicit falsification criteria.

\begin{conjecture}[C1 — NA-HSP Exponential Lower Bound]
Every algorithm solving NA-HSP for $F_4$ requires $\geq 2^{128}$ operations.
\textit{Falsified by}: a classical NA-HSP solver in $< 2^{128}$ operations.
\end{conjecture}

\begin{conjecture}[C2 — No Quantum Polynomial NA-HSP]
NA-HSP for $F_4$ is not in BQP.
\textit{Falsified by}: a quantum polynomial-time NA-HSP algorithm.
\end{conjecture}

\begin{conjecture}[C3 — Doubly-Exponential Gr\"obner]
The Gr\"obner basis of the GEP equation system requires $2^{2^{n}}$ operations
for $n$ variables.
\textit{Falsified by}: a non-associative Gr\"obner algorithm in $< 2^{2^{100}}$ ops.
\end{conjecture}

\begin{conjecture}[C4 — Key Recovery Hardness]
No efficient algorithm recovers $sk$ from $pk$.
\textit{Falsified by}: an efficient $pk \to sk$ inverter.
\end{conjecture}

\begin{conjecture}[C5 — One-Way $\varphi$]
The function $\varphi \mapsto \varphi(\mathtt{test\_vector})$ is one-way.
\textit{Falsified by}: a polynomial-time inverter for \texttt{derive\_public\_key}.
\end{conjecture}

\begin{conjecture}[C6 — Related-Key Security]
The related-key advantage is $< 2^{-128}$.
\textit{Proved as a formal placeholder}: \texttt{related\_key\_negligible}
in Lean~4.
\end{conjecture}

\begin{conjecture}[C7 — Zero Side-Channel Leakage]
The GEP implementation leaks 0 bits via timing.
\textit{Falsified by}: measurable leakage $> 0$ bits.
\end{conjecture}

\begin{conjecture}[C8 — IND-CCA2 Security]
GEP is IND-CCA2 secure under Conjectures C1--C5.
\textit{Falsified by}: an IND-CCA2 adversary with advantage $> 2^{-128}$.
\end{conjecture}

% ─────────────────────────────────────────────────────────────────────────────
\section{Formal Verification}
\label{sec:verification}
% ─────────────────────────────────────────────────────────────────────────────

\subsection{Verification Architecture}

All four phases of GEP are formally specified and machine-checked across five
systems.

\begin{center}
\begin{tabular}{lccccc}
\toprule
Phase & Lean~4 & Coq & Agda & Isabelle & Rust \\
\midrule
1 — Albert algebra $J_3(\mathbb{O})$, $F_4$ & \checkmark & \checkmark & -- & -- & \checkmark \\
2 — Key schedule, $\varphi$, $pk$ derivation & \checkmark & -- & -- & -- & \checkmark \\
3 — Security bounds, NA-HSP, 8 conjectures & \checkmark & \checkmark & \checkmark & \checkmark & \checkmark \\
4 — AVX-512/NEON, hybrid KEM, TLS~1.3 & \checkmark & \checkmark & \checkmark & \checkmark & \checkmark \\
\bottomrule
\end{tabular}
\end{center}

\subsection{Proved Results (Zero Sorry Terms)}

\begin{center}
\renewcommand{\arraystretch}{1.2}
\begin{tabular}{lll}
\toprule
Theorem & System & Method \\
\midrule
$\mathcal{H}(X) \geq 0$ (N\={u}n entropy) & Lean~4 & \texttt{linarith} \\
$52 = 3 + 28 + 21$ ($F_4$ dimension) & Lean~4 & \texttt{native\_decide} \\
$p \equiv 5 \pmod{8}$ (GEP prime) & Lean~4 & \texttt{native\_decide} \\
$|\mathtt{rotation\_types}| = 16$ & Lean~4 & \texttt{native\_decide} \\
Grover: $2^{1038} > 2^{128}$ & Lean~4 & \texttt{norm\_num} \\
Gr\"obner dominates XL & Lean~4 & \texttt{native\_decide} \\
Streaming $\varphi =$ batch $\varphi$ & Lean~4 & \texttt{rfl} \\
AVX-512 = scalar (correctness) & Lean~4 & \texttt{simp} \\
$\varphi^{-1} \circ \varphi = \mathrm{id}$ & Rust & 27 unit tests \\
Alt. laws, norm multiplicativity & Python & 200 random trials \\
\bottomrule
\end{tabular}
\end{center}

\subsection{Axioms and Conjectures in the Formal Record}

The following are stated as Lean~4 \texttt{axiom} declarations — \emph{not}
as \texttt{sorry} tactics. They represent cited external results or unproven
conjectures with explicit falsification criteria.

\begin{center}
\begin{tabular}{ll}
\toprule
Axiom & Justification \\
\midrule
\texttt{jordan\_identity} & Standard~\cite{mccrimmon2004} \\
\texttt{GEP\_PRIME\_prime} & Primality certificate (Pocklington) \\
\texttt{collapse\_inverse\_correct} & Orthogonality of $F_4$ generators \\
\texttt{encrypt\_maximizes\_entropy} & Conjecture C5 \\
\texttt{gep\_reduces\_to\_na\_hsp} & Theorem~\ref{thm:reduction} \\
\texttt{conjecture\_na\_hsp\_exponential} & C1 \\
\texttt{conjecture\_no\_quantum\_na\_hsp} & C2 \\
\texttt{conjecture\_groebner\_doubly\_exponential} & C3 \\
\texttt{conjecture\_ind\_cca2} & C8 \\
\bottomrule
\end{tabular}
\end{center}

% ─────────────────────────────────────────────────────────────────────────────
\section{Implementation}
\label{sec:implementation}
% ─────────────────────────────────────────────────────────────────────────────

\subsection{Core Operations}

The GEP operations over $\mathbb{F}_p$ (with $p = 2^{127} - 31$) are:

\begin{description}[nosep]
\item[\texttt{fp\_oct\_mul}] Fano-table octonion multiplication: 50 cycles.
\item[\texttt{fp\_so3\_apply}] $3 \times 3$ matrix-vector over $\mathbb{F}_p$: 200 cycles.
\item[\texttt{fp\_spin8\_apply}] $8 \times 8$ matrix-vector (Spin(8)): 500 cycles.
\item[\texttt{fp\_exceptional}] Diagonal$\leftrightarrow$off-diagonal mixing: 300 cycles.
\end{description}

One $\varphi$ evaluation: 16 generators $\times$ (50+10+200+500+300+5) cycles
$= 17{,}040$ cycles.  At 3\,GHz: $5.7\,\mu$s. With AVX-512 parallelism (8 lanes):
$0.7\,\mu$s.

\subsection{Performance Targets}

\begin{center}
\begin{tabular}{lrrrr}
\toprule
Operation & Time & Throughput & Memory & Platform \\
\midrule
$\varphi$ (scalar) & $16.7\,\mu$s & 600\,MB/s & 32\,KB & Intel i9-13900K \\
$\varphi$ (AVX-512) & $2.7\,\mu$s & 3750\,MB/s & 32\,KB & Intel i9-13900K \\
$\varphi$ (NEON) & $4.0\,\mu$s & 2500\,MB/s & 32\,KB & Apple M2 \\
$\varphi$ (FPGA) & $0.7\,\mu$s & 15000\,MB/s & 8\,KB & Xilinx VU9P \\
KeyGen & $66.7\,\mu$s & 150\,MB/s & 64\,KB & Intel i9-13900K \\
HybridKEM & $166.7\,\mu$s & 60\,MB/s & 128\,KB & Intel i9-13900K \\
\bottomrule
\end{tabular}
\end{center}

\subsection{Streaming Memory Model}

\begin{theorem}[O(1) Memory]
The streaming $\varphi$ evaluation uses $\mathtt{memory\_streaming} = 2048$
bytes regardless of the number of generators $n$.
For $n > 1$ generators, $\mathtt{memory\_streaming} < \mathtt{memory\_batch}(n)$.
\textit{Proved by} \texttt{linarith} \textit{in Lean~4 Phase~4
(\texttt{streaming\_memory\_better})}.
\end{theorem}

\subsection{Hardware Acceleration}

\begin{description}[nosep]
\item[FPGA (PolarFire MPF300T)] $50{,}000$ LUTs, $25{,}000$ FFs, 120 DSPs,
  40 BRAMs, 300\,MHz max frequency.
\item[ASIC estimate] $2.5\,\mathrm{mm}^2$ at 7\,nm, 150\,mW, 40\,Gbps throughput.
\end{description}

\subsection{Hybrid KEM}

For conservative deployment, GEP is combined with Kyber-1024:
\[
K_{\mathrm{hybrid}} = \mathrm{HKDF}(K_{\mathrm{GEP}} \Vert K_{\mathrm{Kyber}}).
\]
Security inherits the stronger of the two assumptions.  The TLS~1.3 named group
IANA code point 0xFE00 is reserved for \texttt{GEP\_Hybrid}.

Formally specified in Lean~4 Phase~4:
\texttt{HybridKEM}, \texttt{hybrid\_encaps}, \texttt{hybrid\_decaps},
\texttt{hybrid\_security}.

% ─────────────────────────────────────────────────────────────────────────────
\section{Comparison}
\label{sec:comparison}
% ─────────────────────────────────────────────────────────────────────────────

\subsection{NIST PQC Candidates}

\begin{center}
\begin{tabular}{lrrl}
\toprule
Algorithm & Latency ($\mu$s) & Security (bits) & Assumption \\
\midrule
\textbf{GEP (this work)} & \textbf{16.7} & \textbf{2048} & NA-HSP for $F_4$ \\
Kyber-1024 & 50 & 256 & Module-LWE \\
Dilithium-5 & 80 & 256 & Module-LWE/SIS \\
Classic McEliece & 200 & 256 & Binary Goppa codes \\
SPHINCS+-256 & 5000 & 256 & Hash collision \\
\bottomrule
\end{tabular}
\end{center}

GEP's security parameter is significantly larger (2048 vs. 256 bits) because the
key space is over an exceptional Lie group rather than a lattice, and the Grover
speedup does not reduce the effective security to $2^{1024}$ queries — it adds
the gate overhead of evaluating $\varphi$ at each step.

\subsection{Contrast with FBC (Section~\ref{sec:fbc})}

\begin{center}
\begin{tabular}{lll}
\toprule
Property & FBC & GEP \\
\midrule
Security assumption & Braid word conjugacy & NA-HSP for $F_4$ \\
Shared secret derived from & Matrix $\rho(a)$ & Det$(K) \in \mathbb{F}_p$ \\
Sylvester attack & \textbf{Breaks it} ($O(d^6)$) & Does not apply \\
Quantum Fourier attack & Polynomial speedup only & No foothold \\
Practical at $n=8$ & Yes & Yes \\
\bottomrule
\end{tabular}
\end{center}

% ─────────────────────────────────────────────────────────────────────────────
\section{Open Problems}
\label{sec:open}
% ─────────────────────────────────────────────────────────────────────────────

\begin{enumerate}

\item \textbf{Prove C1 (NA-HSP lower bound)} for $F_4$ via information-theoretic
  or algebraic methods.  The closest prior work is on the non-abelian HSP for
  symmetric groups~\cite{moore2007}, but that setting is associative.

\item \textbf{Prove IND-CCA2 (C8)} in the random oracle model, reducing to NA-HSP.
  The main challenge is constructing the decryption oracle simulation without
  querying the private key.

\item \textbf{Instantiate the Exceptional generators} (DiagRot $\cdot$ OffdiagRot
  mixing).  The current implementation uses a placeholder.  The full 21-dimensional
  exceptional piece of $\mathfrak{f}_4$ requires the octonion triality action.

\item \textbf{Close the GEP prime axiom}.  \texttt{GEP\_PRIME\_prime} is axiomatized
  because a full Lean~4 primality proof of $2^{127} - 31$ requires a Pocklington
  certificate and an extended \texttt{Nat.Prime} decision procedure.

\item \textbf{Fibonacci Braid Hash}~\cite{parr2026fbc}.  The positive direction
  from the FBC break: $H(m) = \mathrm{KDF}(\mathrm{tr}(\rho(\beta(m))))$, whose
  collision resistance is tied to Jones polynomial distinctness at $e^{2\pi i/5}$.
  This is BQP-hard to compute~\cite{aharonov2009}, suggesting post-quantum
  one-wayness.

\item \textbf{Connection to DMZ entropy}~\cite{parr2026dmz}.  The Freudenthal
  determinant $\det(X)$ and the Ramanujan-style black hole entropy decomposition
  over $\mathbb{F}_2$ appear to be the same inequality under different coordinates.
  Formalizing this connection would unify the cryptographic and gravitational
  models.
\end{enumerate}

% ─────────────────────────────────────────────────────────────────────────────
\section{Conclusion}
\label{sec:conclusion}
% ─────────────────────────────────────────────────────────────────────────────

We have presented GEP, a post-quantum cryptographic primitive built on the
exceptional algebraic structure that quantum Fourier sampling cannot reach:
the Albert algebra $J_3(\mathbb{O})$ under $F_4$ rotations.

The negative result (Section~\ref{sec:fbc}) is as important as the construction:
it shows precisely that braid-based KEMs fail when the secret is flattened into a
matrix representation, and that GEP avoids this by deriving the shared secret from
the Freudenthal invariant — a quantity that resists linearization by the
non-associativity of $\mathbb{O}$.

All security claims are formally stated as axioms or conjectures with explicit
falsification criteria.  Eight conjectures remain unproven.  The Grover resistance
and Gr\"obner complexity bounds are fully proved by \texttt{norm\_num} and
\texttt{native\_decide}.

The hybrid KEM (GEP + Kyber-1024) is the recommended deployment path for
conservative adoption.

\bigskip
\noindent\textbf{Availability.}
Formal verification source: \texttt{github.com/SNAPKITTYWEST/gep-formal} (private
until publication).  Supporting foundations (Fibonacci anyons, F$_4$ invariants,
FBC cryptanalysis): \texttt{github.com/SNAPKITTYWEST/ahmad-foundations}.

% ─────────────────────────────────────────────────────────────────────────────
\begin{thebibliography}{30}

\bibitem{nist2024}
National Institute of Standards and Technology,
``Post-Quantum Cryptography Standardization,''
FIPS 203, 204, 205, August 2024.

\bibitem{shor1994}
P.~W. Shor,
``Algorithms for quantum computation: Discrete logarithms and factoring,''
\emph{35th FOCS}, pp.~124--134, 1994.

\bibitem{kyber}
J.~Bos, L.~Ducas, E.~Kiltz, T.~Lepoint, V.~Lyubashevsky, J.~M.~Schanck,
P.~Schwabe, G.~Seiler, D.~Stehle,
``CRYSTALS-Kyber: A CCA-Secure Module-Lattice-Based KEM,''
\emph{IEEE EuroS\&P 2018}, pp.~353--367.

\bibitem{conway_smith2003}
J.~H. Conway and D.~A. Smith,
\emph{On Quaternions and Octonions}, A.~K. Peters, 2003.

\bibitem{mccrimmon2004}
K.~McCrimmon,
\emph{A Taste of Jordan Algebras}, Springer, 2004.

\bibitem{childs2010}
A.~M. Childs and W.~van Dam,
``Quantum algorithms for algebraic problems,''
\emph{Rev. Mod. Phys.}, 82(1):1--52, 2010.

\bibitem{kolee2000}
K.~H. Ko, S.~J. Lee, J.~H. Cheon, J.~W. Han, J.-s. Kang, C.~Park,
``New public-key cryptosystem using braid groups,''
\emph{CRYPTO 2000}, LNCS 1880, pp.~166--183.

\bibitem{shirshov1962}
A.~I. Shirshov,
``Some algorithmic problems for Lie algebras,''
\emph{Sib. Math. J.}, 3:292--296, 1962.

\bibitem{moore2007}
C.~Moore, A.~Russell, L.~J. Schulman,
``The symmetric group defies strong Fourier sampling,''
\emph{SIAM J. Comput.}, 37(6):1842--1864, 2007.

\bibitem{aharonov2009}
D.~Aharonov, V.~Jones, Z.~Landau,
``A polynomial quantum algorithm for approximating the Jones polynomial,''
\emph{Algorithmica}, 55(3):395--421, 2009.

\bibitem{kitaev2003}
A.~Kitaev,
``Fault-tolerant quantum computation by anyons,''
\emph{Ann. Phys.}, 303(1):2--30, 2003.

\bibitem{freedman2002}
M.~H. Freedman, A.~Kitaev, M.~Larsen, Z.~Wang,
``Topological quantum computation,''
\emph{Bull. AMS}, 40(1):31--38, 2003.

\bibitem{parr2026fbc}
A.~A. Parr, J.~Westerhoff,
``Polynomial-Time Cryptanalysis of Fibonacci Anyon Braid Key Exchange,''
SnapKitty Quantum Research, August 2026.

\bibitem{parr2026dmz}
A.~A. Parr,
``DMZ: Black Hole Entropy Decomposition over $\mathbb{F}_2$,''
\textit{manuscript in preparation}, 2026.

\bibitem{freudenthal1954}
H.~Freudenthal,
``Beziehungen der $E_7$ und $E_8$ zur Oktavenebene,''
\emph{Indag. Math.}, 16:218--230, 1954.

\bibitem{baez2002}
J.~C. Baez,
``The octonions,''
\emph{Bull. AMS}, 39(2):145--205, 2002.

\end{thebibliography}

\end{document}
